% !TEX root = ../main.tex
\section{Contributions}
\label{sec:contributions}

The choice and organization of the latent representation are crucial design decisions for continuous diffusion reasoning models. We develop this perspective through a complete reasoning pipeline, connecting representation learning, denoising dynamics, and the timing of prompt adaptation. Our main contributions are:

\begin{enumerate}
    \item \textbf{A complete recipe for competitive continuous-diffusion reasoning.}
    We present a practical training and inference recipe for mathematical reasoning and code generation. Our models exceed reported results of recent continuous diffusion models at comparable denoising-backbone scales and achieve competitive performance with selected substantially larger continuous and discrete diffusion models. Strong supervised results, followed by further gains from reward-based post-training, provide evidence that continuous denoising can support challenging reasoning tasks.
    \emph{[Figure~\ref{fig:pipelineoverview}: pipeline overview; Section~\ref{sec:methodology}: training recipe; Section~\ref{sec:headlineresults} and Tables~\ref{tab:headlinecomparison}--\ref{tab:headlinelownfe}: mathematical reasoning; Table~Y: coding results.]}

    \item \textbf{Representation design enables structured denoising dynamics.}
    We show that the space in which denoising occurs is central to reasoning performance. We learn compact, whitened representations from multiple layers of an autoregressive teacher, preserving information relevant to its token predictions. This layer-wise decomposition enables an asynchronous denoising framework in which different representation groups evolve at different rates. We investigate the roles of these spaces through representation comparisons and layer-wise perturbation analyses, connecting the structure of the latent representation to its noise sensitivity and the design of denoising trajectories.
    \emph{[Section~\ref{sec:representation}: learned representations; Figure~\ref{fig:representationcomparison}: design comparisons; Figure~\ref{fig:representationcka}: layer structure; Section~Y: asynchronous denoising; Figure~Y: noise sensitivity.]}

    \item \textbf{Staged learning of conditional representations.}
    We develop a curriculum that first establishes the flow model under stable teacher conditioning, then adapts a learned prompt encoder jointly with the established model. This removes the teacher Transformer from inference and addresses the coupled problem of learning a conditional generator and its conditioning representation. Prompt imitation, substitution, and subsequent adaptation reveal the importance of when conditioning becomes learnable, motivating a progression from teacher-provided structure to representations adapted to the learned flow. Our coding extension further investigates this interface through a prompt encoder that reuses a pretrained diffusion backbone.
    \emph{[Sections~\ref{sec:conditionalencoding}--\ref{sec:conditioningcurriculum} and Lemma~\ref{lem:conditionalrepresentation}: conditional representation learning; Figure~\ref{fig:promptcurriculum}: training stages and adaptation; Figure~\ref{fig:representationcomparison}: joint training from initialization; Figure~Z: coding encoder comparisons.]}

    \item \textbf{Diffusion RL for continuous reasoning.}
    We adapt DiffusionNFT to learned self-conditioning classifier-free guidance (SCCFG), optimizing corrected vector fields while retaining guided rollouts. Gold-solution anchors and verifiable rewards support joint refinement of the flow model and its prompt encoder; a field-level identity characterizes the anchoring effect of gold samples. Improvements over supervised mathematical-reasoning checkpoints provide an encouraging first step toward effective diffusion-RL post-training for continuous reasoning models. We further examine how these gains relate to changes in answer diversity and coverage across multiple samples.
    \emph{[Section~\ref{sec:nft}: SCCFG-compatible NFT; Lemma~\ref{lem:goldanchor}: gold anchoring; Table~\ref{tab:headlinecomparison}: post-training gains; Figure~\ref{fig:headlinestages}: accuracy and coverage.]}
\end{enumerate}
